\documentclass[11pt]{article}
\usepackage[margin=1in]{geometry}
\usepackage{amsmath,amssymb,booktabs}
\usepackage{hyperref}
\title{EXIM AI Co-Pilot: A Multimodal Cargo Scanning and Customs Automation Framework}
\author{Syed Hafiz Ali}
\date{Research and system-design manuscript in preparation}
\begin{document}
\maketitle
\begin{abstract}
This manuscript describes a proposed evidence-centered architecture for assisting export-import documentation and cargo discrepancy review. It connects document extraction, declaration comparison, cargo observations, provenance, role-based review, and human escalation. The framework is a design and research direction; it does not claim deployed customs integration, certified scanning hardware, regulatory approval, or validated operational performance.
\end{abstract}
\section{Operational motivation}
Trade and logistics workflows combine heterogeneous documents and physical observations. Errors may arise from extraction, inconsistent declarations, stale records, missing evidence, or mismatched cargo attributes. A useful assistant must preserve source provenance and distinguish uncertainty from verified discrepancy.
\section{Proposed architecture}
The framework separates ingestion, document normalization, entity extraction, cross-document consistency checks, cargo-observation comparison, evidence scoring, role-based authorization, and review queues. Each finding should retain source references, timestamps, confidence, and a human-readable reason.
\section{Evidence model}
Let each claim be represented by a value, source, timestamp, extraction method, and review state. Agreement among documents is not necessarily independent corroboration: copied or derived documents may share a common upstream error. The design therefore records source lineage and flags correlated evidence.
\section{Cargo observation boundary}
The software framework may accept externally produced observations such as category, dimensions, weight, seal state, or scanner-derived signals. It does not itself create X-ray, gamma, terahertz, RFID, or other physical sensing capability. Hardware integration and calibration are separate validation requirements.
\section{Evaluation design}
Use synthetic invoices, packing lists, declarations, and cargo observations with labelled missing fields, extraction errors, stale authorizations, copied-source conflicts, and adversarial inconsistencies. Evaluate field extraction separately from discrepancy detection and decision quality. Report precision, recall, calibration, false escalation, missed discrepancy, and evidence traceability.
\section{Governance and human review}
The system should recommend rather than silently execute consequential customs or financial actions. Low-confidence, conflicting, stale, or high-impact cases should be routed to authorized human reviewers with a complete evidence packet and audit record.
\section{Results}
No real-world accuracy, throughput, or deployment results are claimed in this manuscript. Synthetic evaluations must be labelled as such.
\section{Limitations}
The current manuscript is an architecture and evaluation proposal. Access to representative lawful datasets, customs-system interfaces, calibrated scanning devices, security review, and operational validation remains future work.
\section{Conclusion}
The central contribution sought is not a generic document chatbot, but an evidence-preserving workflow for identifying discrepancies and escalating decisions in regulated trade operations. Its utility must be established through transparent evaluation and real stakeholder review.
\begin{thebibliography}{9}
\bibitem{wto2014tfa} World Trade Organization. Agreement on Trade Facilitation, Article 7, including provisions on risk management and release and clearance of goods. Official legal text: \url{https://www.wto.org/english/docs_e/legal_e/tfa-nov14_e.htm}.
\bibitem{wco2025safe} World Customs Organization. SAFE Framework of Standards to Secure and Facilitate Global Trade, 2025 edition. \url{https://www.wcoomd.org/en/topics/facilitation/instrument-and-tools/frameworks-of-standards/safe_package.aspx}.
\bibitem{wco2025nii} World Customs Organization. Guidelines for the Procurement and Deployment of Scanning / Non-Intrusive Inspection Equipment, June 2025. \url{https://www.wcoomd.org/en/topics/facilitation/instrument-and-tools/frameworks-of-standards/safe_package/nii-guidelines.aspx}.
\bibitem{unctad2021risk} United Nations Conference on Trade and Development. Technical Note on the WTO Trade Facilitation Agreement: Article 7.4, Risk Management. \url{https://unctad.org/system/files/official-document/wto-technical-note_ch07a07-04_en.pdf}.
\end{thebibliography}
\end{document}
